\section{Récit du stage}

\begin{guideline}
Le récit reconstitue le déroulement de votre stage de façon structurée. Il ne s'agit pas d'un journal de bord, mais d'une narration analytique qui montre votre compréhension des enjeux, de vos décisions et de votre rôle. Faites apparaître votre contribution personnelle à chaque étape.
\end{guideline}

\subsection{Planning des missions}

\begin{guideline}
Présentez de manière chronologique l'ensemble des missions auxquelles
vous avez participé, y compris celles en dehors de votre sujet principal :
vie de l'entreprise, activités commerciales, R\&D, veille, missions
annexes... Un planning visuel est obligatoire. Il permet d'apprécier
votre intégration réelle dans la structure et votre implication au-delà
de votre seule mission principale.
\end{guideline}

\instruction{[Insérez ici votre diagramme de Gantt ou planning visuel
équivalent. Vous pouvez l'intégrer comme image avec :
\texttt{\textbackslash includegraphics[width=\textbackslash textwidth]\{media/gantt.png\}}
après l'avoir exporté depuis votre outil de planification (Excel,
GanttProject, Monday, Notion...). Le planning doit couvrir toute la
durée du stage et distinguer clairement les différentes catégories
de missions.]}

\subsection{Cartographie de la mission}

\begin{guideline}
Décrivez les objectifs de votre mission, les acteurs impliqués (internes/externes), les ressources à disposition et les défis que vous anticipiez. Montrez que vous avez compris les enjeux dès le départ. Vous pouvez utiliser un schéma ou une carte des parties prenantes.
\end{guideline}

\instruction{[Décrivez les objectifs de la mission, les parties prenantes (commanditaire, utilisateurs, partenaires...), les ressources disponibles (outils, budget, équipe) et les défis identifiés a priori. Montrez votre compréhension du périmètre dès le départ.]}


\subsection{Analyse du besoin}

\begin{guideline}
Expliquez comment vous avez qualifié le besoin : qui est le client/utilisateur ? Quelle est sa demande ? Quelle est la problématique sous-jacente ? Comment avez-vous recueilli et formalisé ce besoin ? Faites le lien avec vos apprentissages méthodologiques (UX, product management, etc.).
\end{guideline}

\instruction{[Décrivez le client ou l'utilisateur cible, la demande initiale et la problématique réelle identifiée. Expliquez votre démarche de qualification du besoin : entretiens, ateliers, analyse de données, benchmark... Citez les annexes correspondantes (ex. : Annexe A.1 --- Compte rendu d'entretien utilisateur).]}


\subsection{Défense de la solution}

\begin{guideline}
Expliquez la solution que vous avez proposée ou implémentée : quelle démarche avez-vous suivie (définition, prototypage, tests) ? Pourquoi ce choix ? Quelle valeur ajoutée apporte-t-il ? Quelles contraintes avez-vous prises en compte ? Mettez en avant votre raisonnement d'ingénieur.
\end{guideline}

\instruction{[Décrivez la ou les solutions envisagées et justifiez le choix retenu. Présentez la démarche suivie (itérations, prototypes, tests...), les critères de décision, la valeur créée et les contraintes gérées (techniques, budgétaires, organisationnelles). Référencez les preuves en annexe.]}


\subsection{Modèle de production}

\begin{guideline}
Décrivez comment la solution a été produite concrètement : quelle est la chaîne de production ? Comment était organisée l'équipe (rôles, validations, rituels) ? Quel était le rythme de travail ? Quelles contraintes de performance, qualité ou ressources avez-vous gérées ? Un schéma d'architecture peut être bienvenu.
\end{guideline}

\instruction{[Expliquez l'organisation du travail au quotidien : méthode agile/cycle en V, outils de collaboration, processus de validation. Décrivez la chaîne technique ou organisationnelle et votre rôle précis dedans. Incluez un schéma si pertinent (ex. : architecture système, pipeline CI/CD, flux de données).]}


\subsection{Analyse des risques de la solution}

\begin{guideline}
Identifiez les principaux risques liés à votre solution ou à votre mission : risques techniques, organisationnels, éthiques, opérationnels. Pour chaque risque, précisez le type, la probabilité, l'impact, et les mesures d'atténuation mises en place. Adoptez une posture d'ingénieur responsable.
\end{guideline}

\begin{center}
\renewcommand{\arraystretch}{1.5}
\begin{tabularx}{\textwidth}{@{}l l l X@{}}
\toprule
\textbf{Risque identifié} & \textbf{Type} & \textbf{Criticité} & \textbf{Mesure d'atténuation} \\
\midrule
\instruction{Ex.: Dépendance API} & \instruction{Technique} & \instruction{Élevée} & \instruction{Mise en cache local + fallback} \\
\instruction{Ajoutez vos risques...} & \instruction{...} & \instruction{...} & \instruction{...} \\
\bottomrule
\end{tabularx}
\end{center}
